\subsection{Сценарий 2: near-real-time-витрина при потоковой записи}

В сценарии сравнивались форматы ORC, Parquet и Vortex на таблице с первичным ключом
(схема журнала событий из 25 столбцов, таблица~\ref{tab:schema}) при потоковой записи
с захватом изменений и параллельных аналитических запросах. Ниже приведены результаты
представительного запуска (три повтора полного цикла, ограничение скорости фазы
обновления --- порядка 100\,000~строк/с, интервал чекпойнтов --- 240~с);
для пояснения поведения под конкурентным чтением привлекаются данные ещё одного
запуска.

\textbf{Запись.} Пропускная способность фазы начальной загрузки (последовательные
ключи, реальная запись на диск) практически одинакова у всех трёх форматов
(таблица~\ref{tab:write}); различие в пределах разброса. В фазе обновления показатели
также близки, но интерпретируются с оговоркой: при постоянном значении поля
последовательности значительная доля обновлений дедуплицируется в буфере записи и не
доходит до хранилища, поэтому величина отражает скорость генерации и дедупликации в
памяти. Вывод: Vortex не уступает Parquet и ORC по пропускной способности записи ---
вторая гипотеза работы подтверждается.

\begin{table}
    \caption{Пропускная способность записи, тыс. строк в секунду (среднее по повторам)}
    \label{tab:write}
    \begin{tabularx}{\textwidth}{|X|c|c|}\hline
    \textbf{Формат} & \textbf{Начальная загрузка (INIT)} & \textbf{Обновление (UPDATE)$^{*}$} \\ \hline
    ORC     & 1604 & 91 \\ \hline
    Parquet & 1532 & 98 \\ \hline
    Vortex  & 1527 & 91 \\ \hline
    \multicolumn{3}{|p{0.95\textwidth}|}{\footnotesize $^{*}$~значение фазы обновления отражает скорость генерации и дедупликации в буфере записи (постоянное поле последовательности), а не скорость записи на диск.} \\ \hline
    \end{tabularx}
\end{table}

\textbf{Чтение: усреднённые показатели.} Средние по всем запросам, выборкам и
повторам времена приведены в таблице~\ref{tab:read-avg}. В фазе FINAL (запись
остановлена) Vortex --- быстрейший в среднем; в фазе DURING (параллельная запись)
картина менее однозначна: ORC и Vortex близки, Parquet заметно медленнее за счёт
отдельных катастрофических выбросов (см. ниже).

\begin{table}
    \caption{Среднее время выполнения запроса, мс (по всем запросам, выборкам и повторам)}
    \label{tab:read-avg}
    \begin{tabularx}{\textwidth}{|X|c|c|}\hline
    \textbf{Формат} & \textbf{Фаза DURING} & \textbf{Фаза FINAL} \\ \hline
    ORC     & 7594  & 7633 \\ \hline
    Parquet & 9481  & 7767 \\ \hline
    Vortex  & 7712  & \textbf{6736} \\ \hline
    \end{tabularx}
\end{table}

\textbf{Чтение: фаза FINAL по запросам.} Подробные результаты приведены в
таблице~\ref{tab:read-final} (полужирным выделено наименьшее время по строке). Vortex
выигрывает на десяти запросах из тринадцати, проигрывая только многомерной агрегации
(\texttt{AGG\_MULTIDIM}, незначительно), диапазонному фильтру по URL
(\texttt{RANGE\_URL}) и одному из соединений (\texttt{JOIN\_REVENUE\_PROFILE}). Особенно
велик отрыв на запросах, чувствительных к упорядоченности данных и к качеству
кодирования строк-префиксов: \texttt{RANGE\_SESSION} (соединение с диапазонным
фильтром по идентификатору сессии) --- Vortex быстрее ORC примерно на \SI{20}{\percent};
\texttt{FUNNEL} (многошаговый запрос) --- примерно на \SI{24}{\percent};
\texttt{AGG\_FILTERED} (подсчёт с предикатом) --- примерно на \SI{24}{\percent}.

\begin{table}
    \caption{Время выполнения запросов в фазе FINAL, мс (среднее по повторам; полужирным --- наименьшее по строке)}
    \label{tab:read-final}
    \begin{tabularx}{\textwidth}{|X|c|c|c|}\hline
    \textbf{Запрос} & \textbf{ORC} & \textbf{Parquet} & \textbf{Vortex} \\ \hline
    \texttt{POINT}                & 3944 & 4735 & \textbf{3766} \\ \hline
    \texttt{FULLSCAN\_PROJ}        & 6633 & 5866 & \textbf{5819} \\ \hline
    \texttt{RANGE\_SESSION}        & 7191 & 7673 & \textbf{5801} \\ \hline
    \texttt{RANGE\_URL}            & 7392 & \textbf{5891} & 6451 \\ \hline
    \texttt{AGG\_SUM\_BYTES}       & 5807 & 5836 & \textbf{5781} \\ \hline
    \texttt{AGG\_COUNT\_DISTINCT}  & 8614 & 9184 & \textbf{7819} \\ \hline
    \texttt{AGG\_FILTERED}         & 9345 & 9226 & \textbf{7115} \\ \hline
    \texttt{AGG\_GROUP\_COUNTRY}   & 5943 & 6819 & \textbf{5804} \\ \hline
    \texttt{AGG\_MULTIDIM}         & 7410 & \textbf{6681} & 7288 \\ \hline
    \texttt{FUNNEL}                & 11807 & 11505 & \textbf{8929} \\ \hline
    \texttt{JOIN\_USER\_ACTIVITY}  & 8224 & 9518 & \textbf{8145} \\ \hline
    \texttt{JOIN\_SESSION\_LATENCY}& 9557 & 8926 & \textbf{7343} \\ \hline
    \texttt{JOIN\_REVENUE\_PROFILE}& \textbf{7369} & 9117 & 7512 \\ \hline
    \multicolumn{1}{|r|}{\textbf{выигрышей}} & 1 & 2 & \textbf{10} \\ \hline
    \end{tabularx}
\end{table}

\textbf{Чтение: фаза DURING по запросам.} Результаты приведены в
таблице~\ref{tab:read-during}. Здесь Vortex выигрывает на шести запросах, ORC --- на
пяти, Parquet --- на двух; однако ключевое наблюдение --- не число выигрышей, а
\textit{разброс}. ORC меняется между фазами DURING и FINAL крайне незначительно
(среднее 7594 против 7633\,мс); Parquet деградирует существенно (9481 против 7767\,мс)
и ловит катастрофические выбросы (например, \texttt{AGG\_COUNT\_DISTINCT} ---
24\,801~мс в фазе DURING против 9184~мс в FINAL). Vortex деградирует
умеренно. В другом запуске под более тяжёлой конкурентной нагрузкой выбросы наблюдались
и у Vortex: время полного сканирования и соединений в отдельных DURING-выборках
вырастало в разы --- это попадание запроса в момент крупного сброса данных по
чекпойнту (100\,000~строк/с $\times$ 240~с $=$ около 24~млн строк
одним сбросом --- значительный всплеск ввода-вывода в объектное хранилище).

\begin{table}
    \caption{Время выполнения запросов в фазе DURING, мс (среднее по выборкам и повторам; полужирным --- наименьшее по строке)}
    \label{tab:read-during}
    \begin{tabularx}{\textwidth}{|X|c|c|c|}\hline
    \textbf{Запрос} & \textbf{ORC} & \textbf{Parquet} & \textbf{Vortex} \\ \hline
    \texttt{POINT}                & 4009  & \textbf{3712}  & 3718 \\ \hline
    \texttt{FULLSCAN\_PROJ}        & 7127  & \textbf{5776}  & 6093 \\ \hline
    \texttt{RANGE\_SESSION}        & \textbf{7358}  & 13651 & 10035 \\ \hline
    \texttt{RANGE\_URL}            & 6647  & 6786  & \textbf{6161} \\ \hline
    \texttt{AGG\_SUM\_BYTES}       & 5820  & 5711  & \textbf{5695} \\ \hline
    \texttt{AGG\_COUNT\_DISTINCT}  & \textbf{8616}  & 24801 & 12805 \\ \hline
    \texttt{AGG\_FILTERED}         & 9157  & 9933  & \textbf{7371} \\ \hline
    \texttt{AGG\_GROUP\_COUNTRY}   & \textbf{6003}  & 7191  & 6518 \\ \hline
    \texttt{AGG\_MULTIDIM}         & \textbf{6577}  & 7264  & 6730 \\ \hline
    \texttt{FUNNEL}                & 12384 & 12135 & \textbf{10355} \\ \hline
    \texttt{JOIN\_USER\_ACTIVITY}  & \textbf{8045}  & 8636  & 8492 \\ \hline
    \texttt{JOIN\_SESSION\_LATENCY}& 8672  & 8902  & \textbf{8011} \\ \hline
    \texttt{JOIN\_REVENUE\_PROFILE}& 8302  & 8757  & \textbf{8277} \\ \hline
    \multicolumn{1}{|r|}{\textbf{выигрышей}} & 5 & 2 & 6 \\ \hline
    \end{tabularx}
\end{table}

\textbf{Промежуточные выводы по сценарию.} Vortex эквивалентен Parquet и ORC по
пропускной способности записи (гипотеза~2 подтверждена); на итоговом снимке Vortex ---
быстрейший формат в среднем и на большинстве запросов, особенно на запросах с
диапазонным фильтром по строке-префиксу и с предикатом перед агрегацией; на чистой
агрегации без селективного предиката Vortex иногда уступает (гипотеза~3 частично
подтверждена и здесь --- \texttt{AGG\_MULTIDIM}); под конкурентной записью все
форматы деградируют, причём ORC --- наименее (зрелый нативный читатель с буферизацией
ввода-вывода), Parquet --- наиболее, с катастрофическими выбросами в моменты сброса
данных по чекпойнту, Vortex --- умеренно (гипотеза~4 подтверждена в части различия
поведения форматов под конкурентным чтением).
